\chapter{Regularización en redes profundas}
\section{Apuntes del Tema 6: Regularización en Redes Profundas}

\subsection{1. El Problema del Sobreajuste (\emph{Overfitting}) en Redes Profundas}

\subsubsection{1.1 Naturaleza del Problema}

\begin{itemize}

\item \textbf{Alta Capacidad Expresiva:} Las redes neuronales profundas poseen millones de parámetros libres. Esta gran capacidad les permite memorizar no solo los patrones subyacentes de los datos, sino también el ruido de entrenamiento.

\end{itemize}

\begin{itemize}

\item **Brecha de Generalización (\emph{Generalization Gap}):** Se define como la diferencia entre la pérdida/error en el conjunto de validación (o prueba) y el conjunto de entrenamiento.

\end{itemize}

\begin{itemize}

\item \emph{Baja varianza / Buen ajuste:} La pérdida de validación disminuye a la par con la de entrenamiento.

\end{itemize}

\begin{itemize}

\item \emph{Alta varianza / Overfitting:} A partir de cierto punto, la pérdida de entrenamiento sigue bajando mientras la pérdida de validación empieza a crecer, ensanchando la brecha de generalización.

\end{itemize}

\begin{itemize}

\item \textbf{Objetivo de la Regularización:} Modificar o restringir el proceso de aprendizaje para reducir el error de prueba/generalización sin incrementar excesivamente el error de entrenamiento.

\end{itemize}

\subsection{2. Decaimiento de Pesos (\emph{Weight Decay} / Regularización \(L_2\) vs. \(L_1\))}

La forma clásica de regularización consiste en añadir un término de penalización \(\Omega(W)\) a la función de pérdida original \(\mathcal{L}(W)\):

\$\$\textbackslash\{\}mathcal\{J\}(W) = \textbackslash\{\}mathcal\{L\}(W) + \textbackslash\{\}lambda \textbackslash\{\}Omega(W)\$\$

donde \(\lambda > 0\) es el hiperparámetro que controla la intensidad de la regularización.

\subsubsection{2.1 Decaimiento de Pesos (\(L_2\) Regularization)}

\begin{itemize}

\item \textbf{Penalización:} Se utiliza la norma \(L_2\) al cuadrado de las matrices de pesos:

\end{itemize}

\$\$\textbackslash\{\}Omega(W) = \textbackslash\{\}frac\{1\}\{2\} \textbackslash\{\}Vert\{\}W\textbackslash\{\}Vert\{\}\_2\textasciicircum{}2 = \textbackslash\{\}frac\{1\}\{2\} \textbackslash\{\}sum\_\{i\} \textbackslash\{\}sum\_\{j\} w\_\{ij\}\textasciicircum{}2\$\$

\begin{itemize}

\item \textbf{Efecto en el Gradiente y Actualización (SGD):}

\end{itemize}

\$\$\textbackslash\{\}frac\{\textbackslash\{\}partial \textbackslash\{\}mathcal\{J\}\}\{\textbackslash\{\}partial W\} = \textbackslash\{\}frac\{\textbackslash\{\}partial \textbackslash\{\}mathcal\{L\}\}\{\textbackslash\{\}partial W\} + \textbackslash\{\}lambda W\$\$

\$\$W \textbackslash\{\}leftarrow W - \textbackslash\{\}eta \textbackslash\{\}frac\{\textbackslash\{\}partial \textbackslash\{\}mathcal\{J\}\}\{\textbackslash\{\}partial W\} = W(1 - \textbackslash\{\}eta \textbackslash\{\}lambda) - \textbackslash\{\}eta \textbackslash\{\}frac\{\textbackslash\{\}partial \textbackslash\{\}mathcal\{L\}\}\{\textbackslash\{\}partial W\}\$\$

\begin{itemize}

\item \textbf{Interpretación:} Antes de aplicar el paso de gradiente habitual, los pesos se reducen por un factor multiplicativo \((1 - \eta \lambda) < 1\). Esto fuerza a la red a mantener pesos pequeños y distribuidos de forma suave, previniendo comportamientos extremos frente a pequeñas variaciones en la entrada.

\end{itemize}

\subsubsection{2.2 Regularización \(L_1\) (Lasso)}

\begin{itemize}

\item \textbf{Penalización:} Suma de los valores absolutos de los pesos:

\end{itemize}

\$\$\textbackslash\{\}Omega(W) = \textbackslash\{\}Vert\{\}W\textbackslash\{\}Vert\{\}\_1 = \textbackslash\{\}sum\_\{i\} \textbackslash\{\}sum\_\{j\} \textbackslash\{\}vert\{\}w\_\{ij\}\textbackslash\{\}vert\{\}\$\$

\begin{itemize}

\item \textbf{Efecto:} Mantiene una magnitud de penalización constante \(\lambda \cdot \text{signo}(W)\) independientemente del tamaño del peso, lo que empuja muchos pesos \textbf{exactamente a cero}, generando matrices de pesos esparsas (\emph{sparse models}) y actuando como selección implícita de características.

\end{itemize}

\subsection{3. Dropout (Inyección de Ruido en Activaciones)}

Proposed por Srivastava et al. (2014), es una de las técnicas de regularización más efectivas en redes profundas.

\subsubsection{3.1 Mecanismo de Funcionamiento}

Durante el entrenamiento, en cada iteración o minilote, cada neurona oculta se "apaga" (se fuerza a cero) con una probabilidad \(p \in (0, 1)\) elegida de forma independiente.

\begin{itemize}

\item \textbf{Formulación Matemática (Inverted Dropout):}

\end{itemize}

Dada una activación intermedia \(h\), se transforma en \(h'\) mediante:

\$\$h' = \textbackslash\{\}begin\{cases\} 0 \& \textbackslash\{\}text\{con probabilidad \} p \textbackslash\{\}\textbackslash\{\} \textbackslash\{\}frac\{h\}\{1-p\} \& \textbackslash\{\}text\{con probabilidad \} 1-p \textbackslash\{\}end\{cases\}\$\$

\begin{itemize}

\item \textbf{Reescalado (\(\frac{1}{1-p}\)):} Escalar por \(\frac{1}{1-p}\) durante el entrenamiento garantiza que el valor esperado de la activación sea constante (\(E[h'] = h\)), evitando modificar la magnitud de las entradas en la fase de test.

\end{itemize}

\begin{lstlisting}[language=Python]

  Antes de Dropout                Después de Dropout (con p = 0.4)

    [x1]  [x2]                      [x1]  [x2]

     │ ╲  ╱ │                        │      │

    [h1]  [h2]                     [ 0 ]  [h2]   <- Neurona h1 desactivada

     │ ╱  ╲ │                        │      │

    [o1]  [o2]                      [o1]  [o2]



\end{lstlisting}

\subsubsection{3.2 ¿Por qué funciona Dropout?}

\begin{enumerate}

\item \textbf{Prevención de la Co-adaptación:} Fuerza a cada neurona a aprender características útiles por sí misma, en lugar de depender de la presencia corregidora de otras neuronas específicas.

\end{enumerate}

\begin{enumerate}

\item **Ensamble Implícito (\emph{Implicit Ensembling}):** Entrenar con Dropout equivale a entrenar en cada paso un subconjunto distinto de la red (\emph{thinned network}) que comparte pesos. La evaluación final sin Dropout equivale geométricamente a promediar las predicciones de un número masivo de subredes.

\end{enumerate}

\begin{enumerate}

\item \textbf{Inferencia / Test:} Durante la fase de evaluación, el Dropout debe \textbf{desactivarse} llamando a \texttt{model.eval()} en PyTorch.

\end{enumerate}

\subsection{4. Parada Temprana (\emph{Early Stopping})}

Es una técnica de regularización no paramétrica que aprovecha la dinámica temporal del entrenamiento.

\subsubsection{4.1 Criterio de Funcionamiento}

\begin{itemize}

\item \textbf{Fundamento:} A medida que aumentan las épocas de entrenamiento, el error de entrenamiento disminuye de manera monótona, pero el error de validación alcanza un mínimo y luego empieza a aumentar.

\end{itemize}

\begin{itemize}

\item \textbf{Mecanismo:}

\end{itemize}

\begin{enumerate}

\item Se evalúa el error sobre el conjunto de validación al final de cada época.

\end{enumerate}

\begin{enumerate}

\item Se guarda una copia de los mejores pesos del modelo cuando la pérdida de validación alcanza un nuevo mínimo.

\item **Criterio de Paciencia (\emph{Patience}):** Si la pérdida de validación no mejora durante un número \(P\) determinado de épocas consecutivas, el entrenamiento se interrumpe y se restauran los pesos óptimos guardados.

\end{enumerate}

\begin{lstlisting}[language=Python]

Pérdida

  │        \         /  <- Pérdida de Validación (Empieza a subir)

  │         \_______/

  │          ▲ Parada Temprana (Guardar mejor modelo)

  │           \

  │            \_______ <- Pérdida de Entrenamiento (Sigue bajando)

  └───────────────────────────── Épocas



\end{lstlisting}

\subsection{5. Práctica y Aplicación en Código (PyTorch)}

A continuación se presenta la implementación modular que integra \textbf{Weight Decay (\(L_2\))}, \textbf{Dropout} y \textbf{Early Stopping} utilizando PyTorch:

\begin{lstlisting}[language=Python]

import torch

import torch.nn as nn

import torch.optim as optim

from torch.utils.data import DataLoader, TensorDataset



# --------------------------------------------------------------------------

# 1. Definición del Modelo con Dropout

# --------------------------------------------------------------------------

class RegularizedMLP(nn.Module):

    def __init__(self, input_dim=20, hidden_dim=64, num_classes=2, dropout_p=0.5):

        super(RegularizedMLP, self).__init__()

        

        self.network = nn.Sequential(

            nn.Linear(input_dim, hidden_dim),

            nn.ReLU(),

            nn.Dropout(p=dropout_p),  # Inyección de ruido en activaciones

            

            nn.Linear(hidden_dim, hidden_dim),

            nn.ReLU(),

            nn.Dropout(p=dropout_p),

            

            nn.Linear(hidden_dim, num_classes)  # Logits de salida

        )



    def forward(self, x):

        return self.network(x)



# --------------------------------------------------------------------------

# 2. Clase para Gestionar Early Stopping

# --------------------------------------------------------------------------

class EarlyStopping:

    def __init__(self, patience=5, delta=0.001):

        self.patience = patience

        self.delta = delta

        self.best_loss = float('inf')

        self.counter = 0

        self.early_stop = False

        self.best_state_dict = None



    def __call__(self, val_loss, model):

        if val_loss < self.best_loss - self.delta:

            self.best_loss = val_loss

            self.counter = 0

            self.best_state_dict = model.state_dict().copy() # Guardar mejores pesos

        else:

            self.counter += 1

            if self.counter >= self.patience:

                self.early_stop = True



# --------------------------------------------------------------------------

# 3. Bucle de Entrenamiento y Validación

# --------------------------------------------------------------------------

def train_with_regularization():

    device = torch.device('cuda' if torch.cuda.is_available() else 'cpu')

    

    # Datos sintéticos

    X_train = torch.randn(800, 20)

    y_train = torch.randint(0, 2, (800,))

    X_val = torch.randn(200, 20)

    y_val = torch.randint(0, 2, (200,))

    

    train_loader = DataLoader(TensorDataset(X_train, y_train), batch_size=32, shuffle=True)

    val_loader = DataLoader(TensorDataset(X_val, y_val), batch_size=32, shuffle=False)

    

    model = RegularizedMLP(input_dim=20, dropout_p=0.3).to(device)

    criterion = nn.CrossEntropyLoss()

    

    # Regularización L2 (Weight Decay) aplicada directamente en el optimizador

    optimizer = optim.Adam(model.parameters(), lr=0.001, weight_decay=1e-4)

    early_stopper = EarlyStopping(patience=10)

    

    epochs = 100

    for epoch in range(1, epochs + 1):

        # FASE DE ENTRENAMIENTO (Dropout Activo)

        model.train()

        train_loss = 0.0

        for images, labels in train_loader:

            images, labels = images.to(device), labels.to(device)

            optimizer.zero_grad()

            outputs = model(images)

            loss = criterion(outputs, labels)

            loss.backward()

            optimizer.step()

            train_loss += loss.item() * images.size(0)

            

        train_loss /= len(train_loader.dataset)

        

        # FASE DE VALIDACIÓN (Dropout Desactivado)

        model.eval()  # Importante: desactiva Dropout

        val_loss = 0.0

        with torch.no_grad():

            for images, labels in val_loader:

                images, labels = images.to(device), labels.to(device)

                outputs = model(images)

                loss = criterion(outputs, labels)

                val_loss += loss.item() * images.size(0)

                

        val_loss /= len(val_loader.dataset)

        

        print(f"Época [{epoch:3d}/{epochs}] | Train Loss: {train_loss:.4f} | Val Loss: {val_loss:.4f}")

        

        # Comprobar criterio de Parada Temprana

        early_stopper(val_loss, model)

        if early_stopper.early_stop:

            print(f"\n[Early Stopping] Entrenamiento detenido en la época {epoch}.")

            # Restaurar mejores pesos

            model.load_state_dict(early_stopper.best_state_dict)

            break



if __name__ == '__main__':

    train_with_regularization()



\end{lstlisting}
\section{Conexión con las prácticas}

La práctica permite relacionar las curvas de entrenamiento/validación con weight decay, dropout y early stopping. La interpretación se limita a las métricas y figuras guardadas en el notebook.

\subsection{Fuente práctica: \texttt{06\_Regularization\_01.ipynb}}

Ruta: code/06\_02/06\_Regularization\_01.ipynb

\textbf{Apartados detectados:} Assignment.

Salidas gráficas almacenadas: 13.

\textbf{Salida textual conservada:} \begin{quote}Training model: 100\%|██████████████████████████████████████████████████████████████████████| 5000/5000 [00:13<00:00, 363.98epoch/s, Train Loss=0.0005, Test Loss=1.5976, Train Acc=100.0\%, Test Acc=79.5\%] | Training completed in 13.74 seconds. | <Figure size 1200x400 with 2 Axes> | <Figure size 1200x400 with 2 Axes>\end{quote}

\textbf{Fragmento representativo:}

\begin{lstlisting}[language=Python]

import torch.nn as nn
import torch
import torch.optim as optim
from tqdm import tqdm
import matplotlib.pyplot as plt
import time
import numpy as np

fashion_mnist = torch.load('fashion_mnist.pt')
X_train = fashion_mnist['X_train']
y_train = fashion_mnist['y_train']
X_test = fashion_mnist['X_test']
y_test = fashion_mnist['y_test']

\end{lstlisting}
\section{Autoevaluación razonada}

\subsection*{Pregunta} ¿Qué síntoma sugiere sobreajuste?

\textbf{Solución.} La pérdida de entrenamiento sigue mejorando mientras la validación empeora o deja de mejorar.

\subsection*{Pregunta} ¿Qué hace early stopping?

\textbf{Solución.} Detiene el entrenamiento cuando la métrica de validación no mejora según un criterio y puede restaurar el mejor estado.
\section{Fuentes y trazabilidad}

\textbf{Diapositivas:}

\begin{itemize}

\item \path{slides/06\_01.pdf}

\end{itemize}

\textbf{Nota personal:} apuntes/Apuntes del Tema 6 - Regularización en Redes Profundas.md